% r03_hyb_setup: proof of \Cref{rot:prop:BD}, part I: the hybrid block-spin operators at a gauge field; lengths of the contour pieces; the block-spin algebra for any contour system.
% Input from sec/r03_fill_K5a.tex after the statement of \Cref{rot:prop:BD}. Parts II-VII follow as sec/r03_hyb_coer.tex, r03_hyb_domdiff.tex, r03_hyb_A0.tex, r03_hyb_walk.tex, r03_hyb_ES.tex, r03_hyb_gain.tex.

\paragraph{The hybrid operators at a gauge field.}
Throughout the proof of \Cref{rot:prop:BD} we work in the unit scale of level $r\ge2$ (\cite[(2.20)]{K1}): the physical spacing $\delta_r=5^r\eps_k$ is the unit, the fine spacing is $\eta'=5^{-r}$, $\eta=5\eta'$, and $\ell_j\defeq5^j\eta'$ ($\ell_r=1$). We write $L\defeq5$ for the block factor where a formula is clearer with it, $d=3$, and $f_\alpha=R^{-1}e_\alpha$ for the frame of $\Yk$. The coarse lattice $\Gamma_r=R^{-1}\Z^3$ is common to both flows; $\Gamma_j\defeq R^{-1}(\ell_j\Z^3)$ for $0\le j\le r$, so that $\Gamma_0=\Lambda_Y=R^{-1}(\eta'\Z^3)$ is the fine lattice of the flow $Y$ and $\Gamma_1=R^{-1}(\eta\Z^3)$ is $Y_{k-1}$ rescaled; $\Lambda_X=\eta'\Z^3$ is the fine lattice of the flow $X$ (\Cref{rot:def:match}). Since $5R^{-1}$ is an integer matrix (\Cref{rot:lem:geom}(a)), $\Gamma_1\subset\Lambda_X\cap\Lambda_Y$.

For $1\le j\le r$ and $z\in\Gamma_j$ the \emph{level-$j$ block} of the flow $F$ is
\[
B^F_j(z)\defeq\Lambda_F\cap\big(z+R^{-1}(-\tfrac{\ell_j}2,\tfrac{\ell_j}2)^3\big).
\]
For $F=Y$ this is a cube of $5^{3j}$ points aligned with $(f_\alpha)$; for $F=X$ it is the digitised rotated cube, with no lattice point on its boundary (\Cref{rot:lem:tile}(i)), and $B^X_1(z)=z+\eta'B_0$ is the tile. Both are nested, $B^F_{j+1}(z)=\bigcup_{s\in\{-2,\dots,2\}^3}B^F_j(z+\ell_jR^{-1}s)$, and for $x\in B^F_r(z)$ we write $x_j=x_j(x)$ for the centre of the level-$j$ block containing $x$ ($x_r=z$). Blocks are centred at the sites of the next coarse lattice (\Cref{rot:conv:C}(i)): the root of every block is its centre, in both flows. The matching $\tau:\Lambda_X\to\Lambda_Y$ of \Cref{rot:def:match} maps the tile $B^X_1(x)$ onto the cube $B^Y_1(x)$ for every $x\in\Gamma_1$, with $|\tau x'-x'|\le c_\tau\eta'$, $c_\tau=\sqrt{13}+2\sqrt3<7.1$; points $x''=\tau x'$ are matched. The composite contours are those of \Cref{rot:def:flows}: for the flow $Y$ King's \cite[(2.12)]{K1}, $\Gamma^{Y,(r)}_{z,x}=\Gamma^Y_{z,x_{r-1}}\cup\dots\cup\Gamma^Y_{x_1,x}$, where $\Gamma^Y_{u,v}$ is the L-path that changes the $f_3$-coordinate first, then $f_2$, then $f_1$; for the flow $X$, $\Gamma^{X,(r)}_{z,x'}=\Gamma^X_{z,x_{r-1}}\cup\dots\cup\Gamma^X_{x_2,x_1}\cup\Gamma^X_{x_1,x'}$, where for $i\ge1$ the level-$i$ piece $\Gamma^X_{x_{i+1},x_i}$ is the L-path $\Gamma^Y_{x_{i+1},x_i}$ with every $\eta$-bond $(u,u+\eta f_\alpha)$, $u\in\Gamma_1$, replaced by the fixed staircase $u+\eta'\gamma_\alpha$ (the reversed bond by the reversed staircase), and $\Gamma^X_{x_1,x'}$ is the fixed in-tile path of at most $5$ bonds. The gauge field along a composite contour is read in King's convention \cite[(2.19)]{K1} (\Cref{rot:conv:C}(ii)).

\emph{Backgrounds.} A background is a real bond field $A$ on $\Lambda_F$, identified with a site vector field by $A_{(x,x+\eta'f_\alpha)}=A_\alpha(x)$ (\cite[p.~652]{K1}); $U(A(\Gamma))\defeq\exp[\mathrm ie_rA(\Gamma)]$ with $A(\Gamma)=\sum_{b\in\Gamma}\eta'A_b$ (oriented) and the running charge $e_r=e\,\delta_r^{1/2}$ (\cite[(3.22)]{K1}). \emph{Regular} means \cite[Def.~3.2]{K1}: $|\partial_\mu A_\nu(x)|\le Ce_r^{\beta-1}$ for all $x,\mu,\nu$ and some $\beta>0$. The small-field class at level $r$ is \cite[(3.2)]{K1}: a coarse field $A_r$ on $\Gamma_r$ with $|A_r|\le p_r(\lmu\delta_r)^{-1}$, $|\partial A_r|\le p_r$, $p_r=p(\delta_r)$, lifted to the flow $F$ by \cite[(2.18)]{K1}, $A^{(r),F}\defeq a_rG^{F,\mathrm{ph}}_rQ_F^*A_r$ (the photon fluctuation propagator and the frame-covariant photon average of \Cref{rot:rem:frameav}), plus the Cameron--Martin shift $za_F$ of \Cref{rot:lem:CM}. Both are regular for $\delta_r\le\delta_0$ (\cite[(3.45)]{K1}; \Cref{rot:lem:CM}).

\emph{The pinned operators.} For a flow $F$ and a background $A$ on $\Lambda_F$,
\[
Q_F(A)\phi(z)\defeq\sum_{x\in B^F_r(z)}\eta'^3\,U(A(\Gamma^{F,(r)}_{z,x}))^{-1}\phi(x)\quad(z\in\Gamma_r),\qquad
M_A\defeq-\Delta_A+m_r^2,\qquad G^F_r(A)\defeq\big(M_A+a_rQ_F(A)^*Q_F(A)\big)^{-1},
\]
with $-\Delta_A=D_A^*D_A$ the covariant Laplacian of the fine lattice, $(D_A\phi)(b)=\eta'^{-1}(U(-A_b)\phi(b_+)-\phi(b_-))$, $m_r=m\delta_r$, $a_r=a(1-5^{-2})(1-5^{-2r})^{-1}$, and the pairings $\ip uv=\sum\eta'^3\bar uv$ on $\Lambda_F$ and $\sum_z\bar uv$ on $\Gamma_r$ (unit block volume, $|B^F_r|\eta'^3=1$). This is \cite[(2.11), (2.13)]{K1} in unit scale; the convention $U^{-1}$ makes $Q_F(A)\phi$ transform with the coarse gauge function (\Cref{rot:hyb:lem:alg}). Put
\[
C_A\defeq M_A^{-1},\qquad u^F(A)\defeq C_AQ_F(A)^*,\qquad A^F(A)\defeq\big(a_r^{-1}+Q_F(A)C_AQ_F(A)^*\big)^{-1},\qquad B^F(A)\defeq A^F(A)^{-1},
\]
\[
H^F(A)\defeq a_rG^F_r(A)Q_F(A)^*=u^F(A)A^F(A),\qquad \Gamma^F(A)\defeq C_A-G^F_r(A)=u^F(A)A^F(A)u^F(A)^*=H^FB^FH^{F*}.
\]
By the Woodbury identity these relations are exact for every $A$. $A^F(A)$ is King's effective Laplacian $\Delta^{(r)}(A)=a_r-a_r^2Q_FG_rQ_F^*$ (\cite[(2.14)]{K1}; indeed $a_r(1-a_rQ_FG_rQ_F^*)=a_r(1+a_rQ_FC_AQ_F^*)^{-1}=A^F(A)$), $H^F$ is the lift operator of \cite[(2.18)]{K1} and \Cref{rot:lem:H}, and $C^{(r)}(A)=(A^F(A)+aL^{-2}Q^{(1)}(A)^*Q^{(1)}(A))^{-1}$ is the next fluctuation covariance (\cite[(2.16)]{K1}), $Q^{(1)}$ the one-step average of the coarse lattice. The photon kernels are three copies of the scalar kernels at $A=0$ and mass $\lmu\delta_r$ (\cite[p.~654]{K1}; \Cref{rot:rem:frameav}, \Cref{rot:prop:fc}): no photon kernel sees a background, so the photon sector of \Cref{rot:prop:BD} is the $A=0$ package of \Cref{rot:prop:fc,rot:lem:H,rot:lem:grad,rot:lem:Gam,rot:lem:slice,rot:cor:unif}.

\begin{lemma}[the pieces of the composite contours]\label{rot:hyb:lem:len}
Let $z\in\Gamma_r$, $x'\in B^X_r(z)$, $x''=\tau x'$, and let $x_i$ ($1\le i\le r$) be the common roots.
\begin{enumerate}[label=\textup{(\roman*)},leftmargin=*]
\item For $i\ge1$ the displacement is $x_i-x_{i+1}=\ell_iR^{-1}s$ with $s\in\{-2,\dots,2\}^3$, and every leg of $\Gamma^Y_{x_{i+1},x_i}$ is a concatenation of $|s_\alpha|5^{i-1}$ $\eta$-bonds between points of $\Gamma_1$. The level-$i$ piece has length $\le6\ell_i$ in the flow $Y$ and $\le\frac75\cdot6\ell_i=8.4\,\ell_i$ in the flow $X$; the level-$0$ pieces have length $\le6\eta'$ ($Y$) and $\le5\eta'$ ($X$). The partial contours $\Gamma^{F,(j+1)}_{x_{j+1},x}$ of \cite[(2.19)]{K1} (the pieces of levels $j,j-1,\dots,0$) have length $\le7.5\,\ell_j$ ($Y$) and $\le10.5\,\ell_j$ ($X$); the composite contours have length $\le1.5$ ($Y$) and $\le2.1$ ($X$).
\item Every level-$i$ piece lies inside the open rotated cube $x_{i+1}+R^{-1}(-\frac{\ell_{i+1}}2,\frac{\ell_{i+1}}2)^3$, hence inside $B^F_{i+1}(x_{i+1})$, and the level-$0$ piece inside its first-step cell; so every piece lies in its block and $\Gamma^{F,(r)}_{z,x}\subset B^F_r(z)$.
\item The two composite contours $\Gamma^{X,(r)}_{z,x'}$ and $\Gamma^{Y,(r)}_{z,x''}$ pass through the same points $x_{r-1},\dots,x_1$ and through the same endpoints of $\eta$-bonds; their level-$i$ pieces are within Hausdorff distance $\eta'$ for $i\ge1$, the level-$0$ pieces lie in the ball of radius $\sqrt{13}\eta'$ about $x_1$, and the closed polygon $\Gamma^{X,(r)}_{z,x'}\cup[x',x'']\cup(\Gamma^{Y,(r)}_{z,x''})^{-1}$ bounds a piecewise planar surface $\Sigma(z,x')$, a union of the $\eta$-bond slivers of \Cref{rot:lem:pair} and one polygon inside that ball, of area $\le1.5\eta'+33\eta'^2$. The sub-polygon formed by the partial contours up to level $j+1$ and the matching segment bounds a surface of area $\le1.5\eta'\ell_{j+1}+33\eta'^2$.
\end{enumerate}
\end{lemma}
\begin{proof}
(i) $x_i=\ell_iR^{-1}m_i$ with $m_i\in\Z^3$ the frame coordinates of the block centre and, by nesting, $m_i-5m_{i+1}=:s\in\{-2,\dots,2\}^3$ (the level-$i$ block of $x$ is one of the $125$ sub-blocks of its level-$(i+1)$ block; there are no ties since no lattice point lies on a block face). The L-path from $x_{i+1}$ to $x_i$ has legs $\ell_is_\alpha f_\alpha=5^{i-1}s_\alpha\cdot\eta f_\alpha$, each a run of $|s_\alpha|5^{i-1}$ $\eta$-bonds starting at a point of $\Gamma_1$ (because $x_{i+1}\in\Gamma_{i+1}\subset\Gamma_1$ and $\eta f_\alpha\in\Gamma_1$). Lengths: $\sum_\alpha|s_\alpha|\ell_i\le6\ell_i$; the staircases $\gamma_1,\gamma_2$ have $7$ fine steps for an $\eta$-bond of length $5\eta'$ and $\gamma_3$ has $5$; the in-tile piece has at most $5$ bonds (\Cref{rot:lem:tile}(ii)) and the level-$0$ $Y$-piece at most $6$. Summing, $\sum_{i=1}^{j}8.4\ell_i+5\eta'\le8.4\cdot\frac54\ell_j(1-5^{-j})+5\eta'<10.5\,\ell_j$ and $6\sum_{i=0}^{j}\ell_i\le7.5\,\ell_j$; for $j=r-1$ this gives $2.1$ and $1.5$ (\Cref{rot:lem:pair}). (ii) The endpoints of every $\eta$-bond of $\Gamma^Y_{x_{i+1},x_i}$ are points of $\Gamma_1$ inside the open cube $x_{i+1}+R^{-1}(-\frac{\ell_{i+1}}2,\frac{\ell_{i+1}}2)^3$; in the frame coordinates $(f_\alpha\cdot(\,\cdot\,-x_{i+1}))_\alpha$ they lie in $\eta\Z^3\cap(-\frac{\ell_{i+1}}2,\frac{\ell_{i+1}}2)^3$, and since $\frac{\ell_{i+1}}2=\frac{5^{i+1}}2\eta'$ is an odd multiple of $\frac{\eta'}2$ while the points are multiples of $5\eta'$, each coordinate is at distance $\ge\frac52\eta'$ from the faces. The segment between two such points keeps that distance, and the staircase deviates from the segment by $\le0.6\eta'<\frac52\eta'$ (\Cref{rot:lem:pair}), so it stays in the open cube; its points lie in $\Lambda_X$, hence in $B^X_{i+1}(x_{i+1})$. The in-tile piece lies in the tile (\Cref{rot:lem:tile}(ii)), which is a sub-block. (iii) The Hausdorff distances and the containment of the level-$0$ pieces are \Cref{rot:lem:pair} ($|x'-x_1|\le\sqrt{13}\eta'$, $|x''-x_1|\le2\sqrt3\eta'$, both level-$0$ pieces monotone inside their cells); the two composite contours pass through the same points $x_{r-1},\dots,x_1$ and through the same $\eta$-bond endpoints, so the closed polygon decomposes into the $\eta$-bond slivers (area $\le5\eta'^2$ each, at most $1.5/\eta=0.3/\eta'$ of them, total $\le1.5\eta'$) and the polygon $\Gamma^X_{x_1,x'}\cup[x',x'']\cup(\Gamma^Y_{x_1,x''})^{-1}$ inside the ball of radius $\sqrt{13}\eta'$, of perimeter $\le(5+7.1+6)\eta'$, which the cone from $x_1$ spans with area $\le\frac12\cdot18.1\cdot\sqrt{13}\,\eta'^2<33\eta'^2$. For the partial contours up to level $j+1$ the $Y$-length is $\le7.5\ell_j=1.5\ell_{j+1}$, so there are $\le1.5\ell_{j+1}/\eta$ $\eta$-bonds and the slivers have total area $\le5\eta'^2\cdot1.5\ell_{j+1}/\eta=1.5\eta'\ell_{j+1}$.
\end{proof}

\begin{lemma}[block-spin algebra for any contour system]\label{rot:hyb:lem:alg}
For every background $A$ on $\Lambda_X$ and every system of contours $\Gamma_{z,x}\subset B^X_r(z)$ from the root $z$ to the point $x$:
\begin{enumerate}[label=\textup{(\roman*)},leftmargin=*]
\item $Q_X(A)Q_X(A)^*=\mathbf 1$, and $P_X(A)\defeq Q_X(A)^*Q_X(A)=\sum_zw_zw_z^\dagger$ with $w_z(x)\defeq U(A(\Gamma_{z,x}))\mathbf 1_{B^X_r(z)}(x)$ is an orthogonal projection ($\|w_z\|=1$, orthogonal supports);
\item for a gauge function $\lambda$ on $\Lambda_X$, $Q_X(A+\partial\lambda)(e^{\mathrm ie_r\lambda}\phi)=e^{\mathrm ie_r\lambda|_{\Gamma_r}}Q_X(A)\phi$, hence $G^X_r(A+\partial\lambda)=e^{\mathrm ie_r\lambda}G^X_r(A)e^{-\mathrm ie_r\lambda}$, and likewise for $H^X$, $A^X$, $B^X$, $\Gamma^X$, $C^{(r)}$ with the coarse phases;
\item for a constant field $A\equiv B_0$, $U(B_0(\Gamma_{z,x}))=e^{\mathrm ie_rB_0\cdot(x-z)}$ for every contour, and $G^X_r(B_0)(x,y)=e^{\mathrm ie_rB_0\cdot x}G^X_r(0)(x,y)e^{-\mathrm ie_rB_0\cdot y}$ on $\Lambda_X$ (on a torus: the twisted periodisation of \Cref{rot:hyb:lem:A0}(c));
\item with the contours of \Cref{rot:def:flows}, $Q_X(A)=Q^{\mathrm{rot}}_{r-1}(A)\circ\hat Q(A)$ under the convention \cite[(2.19)]{K1}, and the composite blocks are unions of tiles (\Cref{rot:lem:tile}(iv)).
\end{enumerate}
\end{lemma}
\begin{proof}
(i) is \Cref{rot:lem:tile}(iii),(iv): $|U|=1$ and $|B^X_r(z)|\eta'^3=1$. (ii) $A(\Gamma_{z,x})\mapsto A(\Gamma_{z,x})+\lambda(x)-\lambda(z)$ for any path from $z$ to $x$, so $U(A(\Gamma_{z,x}))^{-1}e^{\mathrm ie_r\lambda(x)}=e^{\mathrm ie_r\lambda(z)}U(A(\Gamma_{z,x}))^{-1}$, and $-\Delta_{A+\partial\lambda}=e^{\mathrm ie_r\lambda}(-\Delta_A)e^{-\mathrm ie_r\lambda}$. (iii) A constant field is $\partial(B_0\cdot x)$ on $\Lambda_X$. (iv) The level-$j$ pieces are the contours of the steps of the flow $Y$ with the bonds re-drawn; iterating the one-step transformations produces exactly the phases of \cite[(2.19)]{K1} along the composite contour (\Cref{rot:conv:C}(ii); \cite[(3.33)--(3.34)]{Ba1}).
\end{proof}

The pairing of the contours is a property of \Cref{rot:def:flows}, not a further hypothesis; a hybrid transformation whose first-step contour were a fine-axis shortest path from the level-$r$ root across the whole block would be a different transformation, for which the contour phases of the two flows differ at order one (\Cref{rot:lem:flux}).
